\begin{tcolorbox}[colback=red!3!white,colframe=red!50!black,boxrule=0.8pt,title={Korrekturhinweise},fonttitle=\bfseries]
\textbf{Korrekturen aus der Rechenprüfung (30.~Sept.~2026).} Die folgenden Stellen sind im Text korrigiert bzw. vermerkt; jede trägt dort einen datierten Hinweis mit dem vorherigen Wert:
\begin{itemize}
\item Tabelle Kosmologie: $E_H/E_P = 1{,}15\times10^{-39} \to 1{,}15\times10^{-61}$; EM-Geometriefaktor $\sqrt{4\pi/9} = 1,18270 \to 1,18164$; $\xi_\ell = 2E_P/(\lambda_C E_P) \to 2\ell_P/\bar{\lambda}_C$; Vermerk zu $\xi_E = 2\sqrt{G}\,E$ mit Werten $10^5$--$10^9$ (in natürlichen Einheiten $\le 2$).
\item (1.~Okt.~2026) Higgs-Ausdruck $\lambda_h^2 v^2/(16\pi^3 E_h^2)$: $1,316$ bzw.\ $1,32\times10^{-4} \to 1,30\times10^{-4}$ (mit $\lambda_h = m_h^2/(2v^2) \approx 0,129$); der Vergleichswert \glqq CODATA/Experimentell\grqq{} $1,320\times10^{-4}$ für $\xi$ existiert nicht und ist durch die Setzung $4/30000$ ersetzt; Übereinstimmung $99,7\,\% \to 97,6\,\%$; \glqq $= \xi$\grqq{} $\to$ \glqq $\approx \xi$\grqq{} (Konsistenzprüfung, Dok.~385).
\end{itemize}
\end{tcolorbox}

\section{Einleitung: Verhältnisbasierte vs. parameterbasierte Physik}

	\medskip\noindent\textbf{Hinweis zur Fassung.} Die deutsche und die englische Fassung dieses Dokuments weichen in Aufbau und Umfang voneinander ab; bei früheren Überarbeitungen wurden in beiden Fassungen Teile gekürzt, ergänzt oder anders gefasst. Bei Abweichungen gilt die umfangreichere deutsche Fassung; Abschnitte, die nur die englische Fassung enthält, gelten als Ergänzung.

	
	Dieses Dokument stellt eine Verifikation des T0-Modells vor, ausgehend von der Auffassung, dass $\xi$ ein Skalenverhältnis ist, kein zugewiesener numerischer Wert. Diese Unterscheidung ist wichtig für das Verständnis der Parametrisierung des T0-Modells, das mit dem einen Parameter $\xi$ auskommt (plus ganzzahlige Strukturwahlen als motivierte Setzungen; SI-Anker dienen der Umrechnung).
	
	\medskip\noindent\textbf{Statusmarker.} Aussagen mit \textbf{[S]} sind \emph{spekulativ}: ein Hinweis oder eine Vermutung, die im Korpus noch nicht hergeleitet oder numerisch geprüft ist.
	
	\begin{tcolorbox}[colback=red!5!white,colframe=red!75!black,title=Zwei Formulierungen]
		\textbf{Übliche Praxis (Literatur):}
		\begin{align}
			\xi &= 1,32 \times 10^{-4} \quad \text{(numerischer Wert zugewiesen)} \\
			\alpha_{\text{EM}} &= \frac{1}{137} \quad \text{(numerischer Wert zugewiesen)} \\
			G &= 6,67 \times 10^{-11} \quad \text{(numerischer Wert zugewiesen)}
		\end{align}
		
		\textbf{T0-Formulierung:}
		\begin{align}
			\xi &\approx \frac{\lambda_h^2 v^2}{16\pi^3 E_h^2} \quad \text{(Higgs-Energieskalenverhältnis, Konsistenzprüfung)} \\
			\xi &= \frac{2\ell_P}{\lambda_C} \quad \text{(Planck-Compton-Längenverhältnis)}
		\end{align}
		\textit{(Korrigiert am 1.~Okt.~2026: vorher $\xi = \lambda_h^2 v^2/(16\pi^3 E_h^2)$ -- mit $\lambda_h = m_h^2/(2v^2) \approx 0,129$ gibt der Higgs-Ausdruck $\xi \approx 1,30\times10^{-4}$, rund $2,3\,\%$ unter $4/30000$; Konsistenzprüfung, keine exakte Herleitung, Dok.~385.)}
	\end{tcolorbox}
	
	\section{Berechnungsverifikation}
	
	Die folgenden Tabellen vergleichen T0-Berechnungen basierend auf Skalenverhältnissen mit etablierten SI-Referenzwerten. Alle Tabellen sind für das Kindle/Portrait-Format skaliert. Spalte \emph{Status}: ja = stimmt überein, neu = neue Vorhersage ohne Messwert, nein = stimmt nicht überein, offen = Abweichung ungeklärt.
	
	% TABELLE 1: Fundamentalkonstanten und abgeleitete Größen
	\begin{table}[H]
		\centering
		\caption{T0-Modell-Verifikation – Teil 1: Fundamentale \& abgeleitete Konstanten}
		\label{tab:verification_fundamental}
		\scriptsize\setlength{\tabcolsep}{3pt}\adjustbox{max width=\textwidth}{%
			\begin{tabular}{p{3.2cm}p{1.2cm}p{2.8cm}p{2.5cm}p{2.3cm}p{1.0cm}p{0.6cm}}
				\toprule
				\textbf{Physikalische Größe} & \textbf{SI-Einheit} & \textbf{T0-Ver-hältnis-formel} & \textbf{T0-Be-rech-nung} & \textbf{CODATA/Experimentell} & \textbf{Übereinstimmung} & \textbf{Status} \\
				\midrule
				$\xi$ (Flach) & 1 & $\xi \approx \frac{\lambda_h^2 v^2}{16\pi^3 E_h^2}$ & $\mathbf{1,30\times10^{-4}}$ & $4/30000$ (Setzung, kein Messwert) & $\mathbf{97,6\%}$ & ja \\
				$\xi$ (Sphärisch) & 1 & $\xi = \frac{\lambda_h^2 v^2}{24\pi^{5/2} E_h^2}$ & $\mathbf{1,557\times10^{-4}}$ & T0-Herleitung & $\mathbf{N/A}$ & neu \\
				Elektronenmasse & MeV & $m_e = f(\xi, \text{Higgs})$ & $\mathbf{0,511}$ & $0,51099895$ & $\mathbf{99,998\%}$ & ja \\
				Compton-Wellenlänge & m & $\lambda_C = \frac{\hbar}{m_e c}$ & $\mathbf{3,862\times10^{-13}}$ & $3,8615927\times10^{-13}$ & $\mathbf{99,989\%}$ & ja \\
				Planck-Länge & m & $\ell_P$ aus $\xi$-Skalierung & $\mathbf{1,616\times10^{-35}}$ & $1,616255\times10^{-35}$ & $\mathbf{99,984\%}$ & ja \\
				\bottomrule
			\end{tabular}
		}
	\end{table}
	
	\textit{(Korrigiert am 1.~Okt.~2026: Zeile $\xi$ (Flach) vorher T0-Berechnung $1,316\times10^{-4}$, Vergleichswert \glqq CODATA/Experimentell\grqq{} $1,320\times10^{-4}$, Übereinstimmung $99,7\,\%$ -- für $\xi$ gibt es keinen CODATA- oder Messwert, der Vergleichswert ist die Setzung $4/30000$; mit $\lambda_h = m_h^2/(2v^2) = 0,1291$ ($m_h = 125,1$~GeV, $v = 246,22$~GeV) gilt $\xi = m_h^2/(64\pi^3 v^2) = 1,301\times10^{-4}$, also $97,6\,\%$; Konsistenzprüfung, Dok.~385.)}
	
	% TABELLE 2: QED-Korrekturen
	\begin{table}[H]
		\centering
		\caption{T0-Modell-Verifikation – Teil 2: QED-Korrekturen}
		\label{tab:verification_qed}
		\scriptsize\setlength{\tabcolsep}{3pt}\adjustbox{max width=\textwidth}{%
			\begin{tabular}{p{3.2cm}p{1.2cm}p{2.8cm}p{2.5cm}p{2.3cm}p{1.0cm}p{0.6cm}}
				\toprule
				\textbf{Physikalische Größe} & \textbf{SI-Einheit} & \textbf{T0-Ver-hältnis-formel} & \textbf{T0-Be-rech-nung} & \textbf{CODATA/Experimentell} & \textbf{Übereinstimmung} & \textbf{Status} \\
				\midrule
				Vertex-Korrektur & 1 & $\frac{\Delta\Gamma}{\Gamma^{\mu}} = \xi^2$ & $\mathbf{1,742\times10^{-8}}$ & Neu & $\mathbf{N/A}$ & neu \\
				Energieunabh. (1 MeV) & 1 & $f(E/E_P)$ & $\mathbf{1,000}$ & Neu & $\mathbf{N/A}$ & neu \\
				Energieunabh. (100 GeV) & 1 & $f(E/E_P)$ & $\mathbf{1,000}$ & Neu & $\mathbf{N/A}$ & neu \\
				\bottomrule
			\end{tabular}
		}
	\end{table}
	
	% TABELLE 3: Kosmologische Vorhersagen
	\begin{table}[H]
		\centering
		\caption{T0-Modell-Verifikation – Teil 3: Kosmologische Vorhersagen [S]}
		\label{tab:verification_cosmological}
		\scriptsize\setlength{\tabcolsep}{3pt}\adjustbox{max width=\textwidth}{%
			\begin{tabular}{p{3.2cm}p{1.6cm}p{2.8cm}p{2.5cm}p{2.3cm}p{1.0cm}p{0.6cm}}
				\toprule
				\textbf{Physikalische Größe} & \textbf{SI-Einheit} & \textbf{T0-Ver-hältnis-formel} & \textbf{T0-Be-rech-nung} & \textbf{CODATA/Experimentell} & \textbf{Übereinstimmung} & \textbf{Status} \\
				\midrule
				$H_0$ (T0) & km/s/Mpc & siehe Dok.~026 & $\mathbf{66{,}2}$ & $67,4$ (Planck) & $\mathbf{1{,}9\%}$ & ja \\
				$H_0$ vs SH0ES & km/s/Mpc & siehe Dok.~026 & $\mathbf{66{,}2}$ & $74,0$ & $\mathbf{10{,}5\%}^*$ & offen \\
				$H_0$ vs H0LiCOW & km/s/Mpc & siehe Dok.~026 & $\mathbf{66{,}2}$ & $73,3$ & $\mathbf{9{,}7\%}^*$ & offen \\
				Universumsalter & Gyr & statisch: $\infty$ & $\mathbf{\infty}$ & $13,8$ & -- & offen \\
				$H_0$ Energieeinheiten & eV & siehe Dok.~026 & $\mathbf{1{,}41\times10^{-33}}$ & T0-Vorhersage & $\mathbf{N/A}$ & neu \\
				$E_H/E_P$-Verhältnis & 1 & siehe Dok.~026 & $\mathbf{1{,}15\times10^{-61}}$ & Theorie & $\mathbf{N/A}$ & neu \\
				\bottomrule
			\end{tabular}
		}
	\end{table}
	
	\textit{(Korrigiert am 30.~Sept.~2026: $E_H/E_P$ vorher $1{,}15\times10^{-39}$ -- $1{,}41\times10^{-33}~\text{eV} / 1{,}2209\times10^{28}~\text{eV} = 1{,}15\times10^{-61}$.)}
	
	\textbf{[S]} Die kosmologischen Werte in Tabelle~\ref{tab:verification_cosmological} sind spekulativ: Der heutige Korpus klammert den kosmischen Sektor bewusst aus, weil auch das Standardmodell ihn nicht herleitet (P39).
	
	% TABELLE 4: Physikalische Felder und Planck-Strom
	\begin{table}[H]
		\centering
		\caption{T0-Modell-Verifikation – Teil 4: Physikalische Felder \& Planck-Strom}
		\label{tab:verification_fields}
		\scriptsize\setlength{\tabcolsep}{3pt}\adjustbox{max width=\textwidth}{%
			\begin{tabular}{p{3.2cm}p{1.6cm}p{2.8cm}p{2.5cm}p{2.3cm}p{1.0cm}p{0.6cm}}
				\toprule
				\textbf{Physikalische Größe} & \textbf{SI-Einheit} & \textbf{T0-Ver-hältnis-formel} & \textbf{T0-Be-rech-nung} & \textbf{CODATA/Experimentell} & \textbf{Übereinstimmung} & \textbf{Status} \\
				\midrule
				Schwinger E-Feld & V/m & $E_S = \frac{m_e^2 c^3}{e\hbar}$ & $\mathbf{1,32\times10^{18}}$ & $1,32\times10^{18}$ & $\mathbf{100,0\%}$ & ja \\
				Kritisches B-Feld & T & $B_c = \frac{m_e^2 c^2}{e\hbar}$ & $\mathbf{4,41\times10^{9}}$ & $4,41\times10^{9}$ & $\mathbf{100,0\%}$ & ja \\
				Planck E-Feld & V/m & $E_P = \sqrt{\frac{c^7}{4\pi\varepsilon_0\hbar G^2}}$ & $\mathbf{6,45\times10^{61}}$ & $6,45\times10^{61}$ & $\mathbf{100,0\%}$ & ja \\
				Planck B-Feld & T & $B_P = E_P/c$ & $\mathbf{2,15\times10^{53}}$ & $2,15\times10^{53}$ & $\mathbf{100,0\%}$ & ja \\
				Planck-Strom (Std) & A & $I_P = \sqrt{\frac{c^6\varepsilon_0}{G}}$ & $\mathbf{9,81\times10^{24}}$ & $3,479\times10^{25}$ & $\mathbf{28,2\%}$ & nein \\
				Planck-Strom (Korr) & A & $I_P = \sqrt{\frac{4\pi c^6\varepsilon_0}{G}}$ & $\mathbf{3,479\times10^{25}}$ & $3,479\times10^{25}$ & $\mathbf{99,98\%}$ & ja \\
				\bottomrule
			\end{tabular}
		}
	\end{table}
	
	\textit{(Korrigiert am 29.~Sept.~2026: Planck-Felder vorher $E_P = c^4/(4\pi\varepsilon_0 G) = 1,04\times10^{61}$~V/m und $B_P = c^3/(4\pi\varepsilon_0 G) = 3,48\times10^{52}$~T -- beide Formeln dimensional unstimmig; richtig ist die Planck-Feldstärke $E_P = \sqrt{c^7/(4\pi\varepsilon_0\hbar G^2)}$ und $B_P = E_P/c$.)}
	
	\section{SI-Planck-Einheitensystem-Verifikation}
	
	\subsection{Komplexe Formelmethode vs. einfache Energiebeziehungen}
	
	\begin{tcolorbox}[colback=yellow!5!white,colframe=yellow!75!black,title=Beobachtung]
		Einfache Beziehungen liefern numerisch genauere Werte als zusammengesetzte Formeln, weil sich weniger Rundungsfehler ansammeln.
	\end{tcolorbox}
	
	% TABELLE 5: Komplexe Formelmethode
	\begin{table}[H]
		\centering
		\caption{SI-Planck-Einheiten: Komplexe Formelmethode}
		\label{tab:verification_complex}
		\scriptsize\setlength{\tabcolsep}{3pt}\adjustbox{max width=\textwidth}{%
			\begin{tabular}{p{3.2cm}p{1.6cm}p{2.8cm}p{2.5cm}p{2.3cm}p{1.0cm}p{0.6cm}}
				\toprule
				\textbf{Physikalische Größe} & \textbf{SI-Einheit} & \textbf{Planck-Formel} & \textbf{T0-Be-rech-nung} & \textbf{CODATA-Referenz} & \textbf{Übereinstimmung} & \textbf{Status} \\
				\midrule
				Planck-Zeit & s & $t_P = \sqrt{\frac{\hbar G}{c^5}}$ & $\mathbf{5,392\times10^{-44}}$ & $5,391\times10^{-44}$ & $\mathbf{100,016\%}$ & ja \\
				Planck-Länge & m & $\ell_P = \sqrt{\frac{\hbar G}{c^3}}$ & $\mathbf{1,617\times10^{-35}}$ & $1,616\times10^{-35}$ & $\mathbf{100,030\%}$ & ja \\
				Planck-Masse & kg & $m_P = \sqrt{\frac{\hbar c}{G}}$ & $\mathbf{2,177\times10^{-8}}$ & $2,176\times10^{-8}$ & $\mathbf{100,044\%}$ & ja \\
				Planck-Temperatur & K & $T_P = \sqrt{\frac{\hbar c^5}{G k_B^2}}$ & $\mathbf{1,417\times10^{32}}$ & $1,417\times10^{32}$ & $\mathbf{99,988\%}$ & ja \\
				Planck-Strom & A & $I_P = \sqrt{\frac{4\pi c^6 \varepsilon_0}{G}}$ & $\mathbf{3,479\times10^{25}}$ & $3,479\times10^{25}$ & $\mathbf{99,980\%}$ & ja \\
				\bottomrule
			\end{tabular}
		}
	\end{table}
	
	\begin{tcolorbox}[colback=gray!5!white,colframe=gray!75!black,title=Anmerkung zu Rundungsfehlern]
		Komplexe Formeln zeigen 99,98–100,04\% Übereinstimmung aufgrund von Rundungsfehlerakkumulation. Dies ist kein Vorhersagefehler, sondern ein Berechnungsartefakt.
	\end{tcolorbox}
	
	\subsection{Methode der einfachen Energiebeziehungen}
	
	% TABELLE 6: Einfache Energiebeziehungen
	\begin{table}[H]
		\centering
		\caption{Natürliche Einheiten: Methode der einfachen Energiebeziehungen}
		\label{tab:verification_simple}
		\scriptsize\setlength{\tabcolsep}{3pt}\adjustbox{max width=\textwidth}{%
			\begin{tabular}{p{3.0cm}p{1.5cm}p{2.8cm}p{2.8cm}p{2.2cm}p{1.2cm}p{0.6cm}}
				\toprule
				\textbf{Physikalische Größe} & \textbf{Beziehung} & \textbf{Beispiel} & \textbf{Elektronenfall} & \textbf{Numerischer Wert} & \textbf{Übereinstimmung} & \textbf{Status} \\
				\midrule
				\multicolumn{7}{l}{\textbf{DIREKTE ENERGIEIDENTITÄTEN - KEINE RUNDUNGSFEHLER}} \\
				\midrule
				Masse & $E = m$ & Energie = Masse & $0,511$ MeV & Gleicher Wert & $\mathbf{100\%}$ & ja \\
				Temperatur & $E = T$ & Energie = Temperatur & $5,93\times10^9$ K & Direkte Konvertierung & $\mathbf{100\%}$ & ja \\
				Frequenz & $E = \omega$ & Energie = Frequenz & $7,76\times10^{20}$ Hz & Direkte Identität & $\mathbf{100\%}$ & ja \\
				\midrule
				\multicolumn{7}{l}{\textbf{INVERSE ENERGIEBEZIEHUNGEN - EXAKT}} \\
				\midrule
				Länge & $E = 1/L$ & Energie = 1/Länge & $3,862\times10^{-13}$ m & Inverse Beziehung & $\mathbf{100\%}$ & ja \\
				Zeit & $E = 1/T$ & Energie = 1/Zeit & $1,288\times10^{-21}$ s & Inverse Beziehung & $\mathbf{100\%}$ & ja \\
				\midrule
				\multicolumn{7}{l}{\textbf{T0-ENERGIEPARAMETER - REINE VERHÄLTNISSE}} \\
				\midrule
				$\xi$ (Higgs, Flach) & $E_h/E_P$ & Energieverhältnis & $1,30\times10^{-4}$ & Higgs-Prüfung zu $4/30000$ & $\mathbf{97,6\%}$ & ja \\
				$\xi$ (Higgs, Sph) & $E_h/E_P$ & Korrigiertes Verhältnis & $1,557\times10^{-4}$ & T0-Herleitung & $\mathbf{100\%}$ & neu \\
				$\xi$ Geometrisch & $E_\ell/E_P$ & Längen-Energie-Verhältnis & $8,37\times10^{-23}$ & Reine Geometrie & $\mathbf{100\%}$ & ja \\
				EM-Geometriefaktor & Verhältnis & $\sqrt{4\pi/9}$ & $1,18164$ & Mathematisch exakt & $\mathbf{100\%}$ & neu \\
				\midrule
				\multicolumn{7}{l}{\textbf{VOLLSTÄNDIGE SI-EINHEITEN-ENERGIEABDECKUNG - ALLE 7/7 EINHEITEN}} \\
				\midrule
				Elektrischer Strom & $I = E/T$ & Energieflussrate & $[E]$ Dim. & Direkte Energiebeziehung & $\mathbf{100\%}$ & ja \\
				Stoffmenge & $[E^2]$ Dim. & Energiedichteverhältnis & Dimensionale Struktur & SI-definierte $N_A$ & $\mathbf{Def.}$ & neu \\
				Lichtstärke & $[E^3]$ Dim. & Energieflusswahrnehmung & Dimensionale Struktur & SI-definierte 683 lm/W & $\mathbf{Def.}$ & neu \\
				\bottomrule
			\end{tabular}
		}
	\end{table}
	
	\textit{(Korrigiert am 30.~Sept.~2026: EM-Geometriefaktor vorher $1,18270$ -- $\sqrt{4\pi/9} = \sqrt{1,396263} = 1,18164$.)}
	\textit{(Korrigiert am 1.~Okt.~2026: Zeile $\xi$ (Higgs, Flach) vorher $1,316\times10^{-4}$, \glqq Aus Higgs-Physik\grqq{}, $100\,\%$ -- mit $\lambda_h = m_h^2/(2v^2)$ folgt $1,30\times10^{-4}$, $97,6\,\%$ von $4/30000$; Konsistenzprüfung.)}
	
	\begin{tcolorbox}[colback=blue!5!white,colframe=blue!75!black,title=Genauigkeit durch Vereinfachung]
		\textbf{Komplexe Formelmethode (Traditionelle Physik):}
		\begin{itemize}
			\item Verwendet: $\sqrt{\frac{\hbar G}{c^5}}$, mehrere Konstanten, Umrechnungsfaktoren
			\item Ergebnis: 99,98–100,04\% Übereinstimmung (Rundungsfehler akkumulieren)
			\item Problem: Jeder Berechnungsschritt führt kleine Fehler ein
		\end{itemize}
		
		\textbf{Methode der einfachen Energiebeziehungen (T0-Physik):}
		\begin{itemize}
			\item Verwendet: Direkte Identitäten $E = m$, $E = 1/L$, $E = 1/T$
			\item Ergebnis: exakt per Definition (Identitäten in natürlichen Einheiten)
			\item Vorteil: Keine Zwischenberechnungen, keine Fehlerakkumulation
		\end{itemize}
		
		\textbf{EINORDNUNG:}
		Der Unterschied ist rechentechnischer Art: Identitäten in natürlichen Einheiten vermeiden Zwischenschritte und damit Rundungsfehler. Das passt zur T0-Auffassung, dass Energie die grundlegende Größe ist; ein Beweis dafür ist es nicht.
		
		\textbf{FAZIT}: Einfache Beziehungen sind hier numerisch nicht weniger genau.
	\end{tcolorbox}
	
	\section{Die $\xi$-Parameter-Hierarchie}
	
	\subsection{Klarstellung}
	
	\begin{tcolorbox}[colback=red!10!white,colframe=red!75!black,title=Hinweis: Verwendung des Symbols $\xi$ in diesem Dokument]
		\textbf{Mögliches Missverständnis:} Behandlung von $\xi$ als universellen Parameter
		
		\textbf{In diesem Dokument:} $\xi$ ist eine \textbf{Klasse dimensionsloser Skalenverhältnisse}, kein einzelner Wert.
		
		\textbf{Folge einer Verwechslung:} Fehlinterpretationen, falsche Zahlenwerte, Dimensionsfehler.
		
		$\xi$ repräsentiert jedes dimensionslose Verhältnis der Form:
		\begin{equation}
			\xi = \frac{\text{T0-charakteristische Energieskala}}{\text{Referenzenergieskala}}
		\end{equation}
		
		Das T0-Modell verwendet $\xi$, um verschiedene dimensionslose Verhältnisse in unterschiedlichen physikalischen Kontexten zu bezeichnen.
	\end{tcolorbox}
	
	\subsection{Die drei fundamentalen $\xi$-Energieskalen}
	
	\begin{table}[H]
		\centering
		\caption{Die drei fundamentalen $\xi$-Parametertypen im T0-Modell}
		\label{tab:xi_hierarchy}
		\scriptsize\setlength{\tabcolsep}{3pt}\adjustbox{max width=\textwidth}{%
			\begin{tabular}{|p{2.8cm}|p{4.5cm}|p{2.2cm}|p{3.8cm}|}
				\hline
				\textbf{Kontext} & \textbf{Definition} & \textbf{Typischer Wert} & \textbf{Physikalische Bedeutung} \\
				\hline
				\textbf{Energieabhängig} & $\xi_E = 2\sqrt{G} \cdot E$ & $10^5$ bis $10^9$ & Energie-Feld-Kopplung \\
				\hline
				\textbf{Higgs-Sektor} & $\xi_H = \frac{\lambda_h^2 v^2}{16\pi^3 E_h^2}$ & $1,30\times10^{-4}$ & Energieskalenverhältnis \\
				\hline
				\textbf{Skalenhierarchie} & $\xi_\ell = \frac{2\ell_P}{\bar{\lambda}_C}$ & $8,37\times10^{-23}$ & Energiehierarchieverhältnis \\
				\hline
			\end{tabular}
		}
	\end{table}
	
	\textit{(Korrigiert am 1.~Okt.~2026: Higgs-Sektor vorher $1,32\times10^{-4}$ -- mit $\lambda_h = m_h^2/(2v^2) \approx 0,129$ gilt $\xi_H = m_h^2/(64\pi^3 v^2) \approx 1,30\times10^{-4}$, rund $2,3\,\%$ unter $4/30000$.)}
	
	\textit{(Korrigiert am 30.~Sept.~2026: $\xi_\ell$ vorher $2E_P/(\lambda_C E_P)$ -- das kürzt sich zu $2/\lambda_C$ und ist nicht dimensionslos; der Zahlenwert $8,37\times10^{-23}$ ist $2\ell_P/\bar{\lambda}_C = 2\cdot1,616\times10^{-35}~\text{m}/3,862\times10^{-13}~\text{m}$ mit $\bar{\lambda}_C = \hbar/(m_e c)$.)}
	
	\textit{(Vermerk vom 30.~Sept.~2026: $\xi_E = 2\sqrt{G}\cdot E$ mit \enquote{typischem Wert $10^5$ bis $10^9$} ist nicht haltbar -- in natürlichen Einheiten ist $\sqrt{G} = 1/E_P$, also $\xi_E = 2E/E_P \le 2$ für alle $E \le E_P$; Werte $10^5$--$10^9$ sind damit ausgeschlossen, und in SI ist $2\sqrt{G}\,E$ nicht dimensionslos. Das betrifft auch Regel~1 unten.)}
	
	\subsection{Anwendungsregeln}
	
	\begin{tcolorbox}[colback=blue!5!white,colframe=blue!75!black,title=Anwendungsregeln für $\xi$-Parameter (Reine Energie)]
		\textbf{Regel 1: Universelle energieabhängige Systeme (EMPFOHLEN)}
		\begin{equation}
			\text{Verwende } \xi_E = 2\sqrt{G} \cdot E \text{ wobei } E \text{ die relevante Energie ist}
		\end{equation}
		
		\textbf{Regel 2: Kosmologische/Kopplungsvereinheitlichung (SPEZIALFÄLLE)}
		\begin{equation}
			\text{Verwende } \xi_H = 1,32 \times 10^{-4} \text{ (Higgs-Energieverhältnis)}
		\end{equation}
		
		\textbf{Regel 3: Reine Energiehierarchieanalyse (THEORETISCH)}
		\begin{equation}
			\text{Verwende } \xi_\ell = 8,37 \times 10^{-23} \text{ (Energieskalenverhältnis)}
		\end{equation}
		
		\textbf{Anmerkung:} In der Praxis gilt Regel 1 für die weitaus meisten T0-Berechnungen aufgrund der ausgeprägten T0-Skalenhierarchie.
	\end{tcolorbox}
	
	\section{Wichtige Erkenntnisse aus der Verifikation}
	
	\subsection{Hauptergebnisse}
	
	\begin{tcolorbox}[colback=green!5!white,colframe=green!75!black,title=Hauptergebnisse der T0-Verifikation]
		\textbf{1. Skalenverhältnis-Validierung:}
		\begin{itemize}
			\item Etablierte Werte: 99,99\% Übereinstimmung mit CODATA
			\item Geometrisches $\xi$-Verhältnis: 100,003\% Übereinstimmung mit Planck-Compton-Berechnung
			\item Dimensionale Konsistenz der geprüften Größen
		\end{itemize}
		
		\textbf{2. Neue testbare Vorhersagen:}
		\begin{itemize}
			\item QED-Vertex-Verhältnisse: $1,74 \times 10^{-8}$ (energieunabhängig)
			\item \textbf{[S]} Kosmologisches $H_0 \approx 66{,}2\,$km/s/Mpc (siehe Dokument~026)
			\item \textbf{[S]} Rotverschiebungsverhältnisse: 40,5\% spektrale Variation
		\end{itemize}
		
		\textbf{3. Gesamtbewertung:}
		\begin{itemize}
			\item Etablierte Werte: 99,99\% Übereinstimmung
			\item Neue Vorhersagen: 14+ testbare Verhältnisse
			\item Dimensionale Konsistenz: für alle Tabellengrößen gegeben
			\item Skalenverhältnis-Basis: in sich konsistent
		\end{itemize}
	\end{tcolorbox}
	
	\subsection{Experimentelle Testbarkeit}
	
	Die verhältnisbasierte Natur des T0-Modells ermöglicht spezifische experimentelle Tests:
	
	\begin{enumerate}
		\item \textbf{Energieskalenunabhängige QED-Korrekturen}:
		\begin{equation}
			\frac{\Delta\Gamma^{\mu}(E_1)}{\Delta\Gamma^{\mu}(E_2)} = 1 \quad \text{für alle } E_1, E_2 \ll E_P
		\end{equation}
		
		\item \textbf{Kosmologische Skalenverhältnisse} \textbf{[S]}:
		\begin{equation}
			\frac{\kappa}{H_0} = \xi \approx \frac{\lambda_h^2 v^2}{16\pi^3 E_h^2}
		\end{equation}
		\textit{(Korrigiert am 1.~Okt.~2026: vorher $\xi = \lambda_h^2 v^2/(16\pi^3 E_h^2)$ -- der Higgs-Ausdruck gibt $\approx 1,30\times10^{-4}$, rund $2,3\,\%$ unter $4/30000$; Konsistenzprüfung, keine Gleichheit.)}
	\end{enumerate}
	
	\section{Schlussfolgerungen}
	
	Die Verifikation stützt die zentrale Aussage des T0-Modells: \textbf{Die fundamentale Physik lässt sich auf Skalenverhältnisse statt auf zugewiesene Parameter gründen}. Der $\xi$-Parameter charakterisiert diese Proportionalitäten und ermöglicht eine Beschreibung physikalischer Phänomene mit dem einen Parameter $\xi$ (plus ganzzahlige Strukturwahlen als motivierte Setzungen; SI-Anker dienen der Umrechnung).
	
	\begin{thebibliography}{9}
		
		\bibitem{pascher_h0_energy_2025}
		Pascher, J. (2025). \textit{Reine Energieformulierung der $H_0$- und $\kappa$-Parameter im T0-Modell-Rahmenwerk}. \\
		\url{https://github.com/jpascher/T0-Time-Mass-Duality/blob/main/2/pdf/xxx_H0_kappa_De.pdf}
		
		\bibitem{pascher_beta_derivation_2025}
		Pascher, J. (2025). \textit{Feldtheoretische Herleitung des $\beta_T$-Parameters in natürlichen Einheiten ($\hbar = c = 1$)}. \\
		\url{https://github.com/jpascher/T0-Time-Mass-Duality/blob/main/2/pdf/093_DerivationVonBeta_De.pdf}
		
		\bibitem{pascher_elimination_mass_2025}
		Pascher, J. (2025). \textit{Elimination der Masse als dimensionaler Platzhalter im T0-Modell: Auf dem Weg zu wahrhaft parameterfreier Physik}. \\
		\url{https://github.com/jpascher/T0-Time-Mass-Duality/blob/main/2/pdf/052_EliminationOfMass_De.pdf}
		
		\bibitem{pascher_mol_candela_2025}
		Pascher, J. (2025). \textit{T0-Modell: Universelle Energiebeziehungen für Mol- und Candela-Einheiten - Vollständige Herleitung aus Energie-Skalierungsprinzipien}. \\
		\url{https://github.com/jpascher/T0-Time-Mass-Duality/blob/main/2/pdf/062_Moll_Candela_De.pdf}
		
	\end{thebibliography}
